\subsection{Exercise 6 - Lab 2}

Yêu cầu: Mục tiêu chính là giảm thiểu độ phức tạp bên trong hàm ngắt timer để tránh lỗi ngắt lồng nhau gây treo hệ thống. Tạo ra một software timer đếm lùi sau mỗi lần ngắt timer xảy ra mỗi 10ms. Sử dụng timer phần mềm này để thay thế và loại bỏ hàm \texttt{HAL\_Delay(1000)} trong hàm \texttt{main()}, nhằm tạo ra một hệ thống xử lý trễ không đồng bộ.

Ý tưởng: Khai báo các biến toàn cục để quản lý timer phần mềm gồm: \texttt{timer\_counter}, \texttt{timer\_flag} và \texttt{TIMER\_CYCLE}. Hiện thực hàm \texttt{setTimer(int duration)} để nạp giá trị đếm bằng cách lấy thời gian mong muốn chia cho chu kỳ. Hàm \texttt{timer\_run()} để lùi biến \texttt{timer\_counter} lớn hơn 0, khi biến đếm về 0 thì cờ bật lên 1.

\begin{lstlisting}[language=C, caption={Phần code hiện thực}, label={lst:ex6_lab2}]
int timer_counter = 0;
int timer_flag = 0;
int TIMER_CYCLE = 10;

void setTimer(int duration){
    timer_counter = duration / TIMER_CYCLE;
    timer_flag = 0;
}

void timer_run(){
    if(timer_counter > 0){
        timer_counter--;
        if(timer_counter == 0) timer_flag = 1;
    }
}

void HAL_TIM_PeriodElapsedCallback(TIM_HandleTypeDef *htim){
    if(htim->Instance == TIM2){
        timer_run();
        HAL_GPIO_WritePin(GPIOA, GPIO_PIN_6|GPIO_PIN_7|GPIO_PIN_8|GPIO_PIN_9, GPIO_PIN_SET);
        update7SEG(index_led);

        if(second_counter == 2 || second_counter == 4) {
            HAL_GPIO_TogglePin(GPIOA, GPIO_PIN_4);
        }

        second++;
        if(second >= 60) {
            second = 0;
            minute++;
        }

        if(minute >= 60) {
            minute = 0;
            hour++;
        }

        if(hour >= 24) {
            hour = 0;
        }

        updateClockBuffer();
    }
}
\end{lstlisting}

\begin{itemize}
    \item \textbf{Report 1:} Nếu xóa dòng \texttt{setTimer(1000);} trước vòng lặp while, biến \texttt{timer\_counter} ban đầu sẽ bằng 0. Cờ sẽ không bao giờ được bật lên 1, dẫn đến LED không bao giờ nhấp nháy.

    \item \textbf{Report 2:} Nếu đổi thành \texttt{setTimer(1);}, thời gian đếm lùi sẽ là $1/10 = 0$. LED sẽ không nhấp nháy.

    \item \textbf{Report 3:} Nếu đổi thành \texttt{setTimer(10);}, LED sẽ nhấp nháy mỗi 10ms. LED sẽ nhấp nháy nhanh, chúng ta vẫn sẽ không nhìn ra được sự nhấp nháy.
\end{itemize}